\documentclass[../principal.tex]{subfiles}
\graphicspath{{\subfix{../imagenes/}}}

\begin{document}

\chapter{Cursos}

Este proyecto nace desde nuestro trabajo conjunto en dos cursos que propusimos en 2024 en la Escuela de Diseño en la Facultad de Arquitectura, Arte y Diseño de la Universidad Diego Portales, para estudiantes de pregrado en tercer y cuarto año. Estos cursos son de tipo taller, y reciben los nombres:

\begin{itemize}
    \item DIS8644 Taller de diseño de máquinas electrónicas
    \item DIS8645 Taller de diseño de máquinas computacionales
\end{itemize}


Las primeras versiones de estos cursos fueron dictadas en 2025, y al momento de escribir estas palabras en 2026 estamos dictando sus segundas iteraciones. El de máquinas electrónicas lo dictamos los primeros semestres (marzo a julio), y el de máquinas computacionales lo dictamos los segundos semestres (agosto a diciembre).


\section{DIS8644 Máquinas electrónicas}

Este curso se hace cargo en cuatro meses de formación de la enseñanza de conceptos maquinales usando circuitos integrados.

Las evaluaciones grupales son 3:

- Proyecto hecho en una placa de pruebas, sin soldadura
- Proyecto hecho en software de diseño de esquemáticos
- Proyecto hecho con carcasa, soldadura, simulación


\section{DIS8645 Máquinas computacionales}

Este curso se hace cargo en cuatro meses de formación de la enseñanza de conceptos maquinales usando código y metáforas de programación orientada a objetos.

\section{Cursos tangenciales}

En nuestras prácticas docentes recientes hemos dictado otros cursos relacionados.

\subsection{DIS09214 Pensamiento computacional}

Curso semestral obligatorio de segundo año de pregrado en la escuela de diseño, dictado en el primer semestre 2026. Fue dictado en múltiples secciones paralelas, entre ellas una dictada por Matías Serrano en conjunto con el ayudante Santiago Gaete, y otra dictada por Aarón Montoya con la ayudante Vania Paredes.

\subsection{DIS9079 Interacciones inalámbricas}

Curso semestral optativo de pregrado en la escuela de diseño, dictado en el primer semestre 2026 por Aarón Montoya en conjunto con el ayudante Mateo Arce. Enseñamos estrategias de uso de computación inalámbrica con microcontroladores.

\newpage

\end{document}

